\documentclass[10pt,a4paper]{amsart}
\usepackage{amsmath,amssymb,amsthm,graphicx}
\usepackage{verbatim}
\newenvironment{reference}{\comment}{\endcomment}
%\newenvironment{reference}{}{}
\newenvironment{slogan}{\comment}{\endcomment}
\newenvironment{history}{\comment}{\endcomment}

% For commutative diagrams we use Xy-pic
\usepackage[all]{xy}

% We use 2cell for 2-commutative diagrams.
\xyoption{2cell}
% Theorem environments.
%
\theoremstyle{plain}
\newtheorem{theorem}[subsection]{Theorem}
\newtheorem{proposition}[subsection]{Proposition}
\newtheorem{lemma}[subsection]{Lemma}

\theoremstyle{definition}
\newtheorem{definition}[subsection]{Definition}
\newtheorem{example}[subsection]{Example}
\newtheorem{exercise}[subsection]{Exercise}
\newtheorem{situation}[subsection]{Situation}

\theoremstyle{remark}
\newtheorem{remark}[subsection]{Remark}
\newtheorem{remarks}[subsection]{Remarks}

\numberwithin{equation}{subsection}

% Macros
%
\def\lim{\mathop{\mathrm{lim}}\nolimits}
\def\colim{\mathop{\mathrm{colim}}\nolimits}
\def\Spec{\mathop{\mathrm{Spec}}}
\def\Hom{\mathop{\mathrm{Hom}}\nolimits}
\def\Ext{\mathop{\mathrm{Ext}}\nolimits}
\def\SheafHom{\mathop{\mathcal{H}\!\mathit{om}}\nolimits}
\def\SheafExt{\mathop{\mathcal{E}\!\mathit{xt}}\nolimits}
\def\Sch{\mathit{Sch}}
\def\Mor{\mathop{\mathrm{Mor}}\nolimits}
\def\Ob{\mathop{\mathrm{Ob}}\nolimits}
\def\Sh{\mathop{\mathit{Sh}}\nolimits}
\def\NL{\mathop{N\!L}\nolimits}
\def\CH{\mathop{\mathrm{CH}}\nolimits}
\def\proetale{{pro\text{-}\acute{e}tale}}
\def\etale{{\acute{e}tale}}
\def\QCoh{\mathit{QCoh}}
\def\Ker{\mathop{\mathrm{Ker}}}
\def\Im{\mathop{\mathrm{Im}}}
\def\Coker{\mathop{\mathrm{Coker}}}
\def\Coim{\mathop{\mathrm{Coim}}}

% Boxtimes
%
\DeclareMathSymbol{\boxtimes}{\mathbin}{AMSa}{"02}

%
% Macros for moduli stacks/spaces
%
\def\QCohstack{\mathcal{QC}\!\mathit{oh}}
\def\Cohstack{\mathcal{C}\!\mathit{oh}}
\def\Spacesstack{\mathcal{S}\!\mathit{paces}}
\def\Quotfunctor{\mathrm{Quot}}
\def\Hilbfunctor{\mathrm{Hilb}}
\def\Curvesstack{\mathcal{C}\!\mathit{urves}}
\def\Polarizedstack{\mathcal{P}\!\mathit{olarized}}
\def\Complexesstack{\mathcal{C}\!\mathit{omplexes}}
% \Pic is the operator that assigns to X its picard group, usage \Pic(X)
% \Picardstack_{X/B} denotes the Picard stack of X over B
% \Picardfunctor_{X/B} denotes the Picard functor of X over B
\def\Pic{\mathop{\mathrm{Pic}}\nolimits}
\def\Picardstack{\mathcal{P}\!\mathit{ic}}
\def\Picardfunctor{\mathrm{Pic}}
\def\Deformationcategory{\mathcal{D}\!\mathit{ef}}
\usepackage{cleveref}
\crefname{originallemma}{original lemma}{original lemmas}
\Crefname{originallemma}{Original lemma}{Original lemmas}
\makeatletter
\newlabel{lemma-consequence-Artin-Rees-bis}{{0G4R}{source}{}{sourcepointer-1}{}}
\newlabel{lemma-consequence-Artin-Rees-bis@cref}{{[originallemma][0][]0G4R}{[1][0][]source}{}{}{}}
\makeatother
\crefname{originalreference}{original reference}{original references}
\Crefname{originalreference}{Original reference}{Original references}
\makeatletter
\newlabel{derived-remark-truncation-distinguished-triangle}{{08J5}{source}{}{sourcepointer-2}{}}
\newlabel{derived-remark-truncation-distinguished-triangle@cref}{{[originalreference][0][]08J5}{[1][0][]source}{}{}{}}
\makeatother
\crefname{originallemma}{original lemma}{original lemmas}
\Crefname{originallemma}{Original lemma}{Original lemmas}
\makeatletter
\newlabel{lemma-deligne-system-2-out-of-3}{{0G4Q}{source}{}{sourcepointer-3}{}}
\newlabel{lemma-deligne-system-2-out-of-3@cref}{{[originallemma][0][]0G4Q}{[1][0][]source}{}{}{}}
\makeatother
\title{The algebraic criterion for being pro-isomorphic to a Deligne system}
\author{The Stacks Project Authors}
\date{stacks-project:duality.tex:local-snapshot-20261012; GNU FDL 1.2 or later}
\begin{document}
\maketitle
\noindent\textit{Editorial note.} All delivered mathematical text is the author's own, copied from the pinned Stacks Project chapter file duality.tex; no statement or proof step was rewritten, reordered or removed.\par\medskip
\noindent\textit{Editorial note.} Preamble provenance: the pinned chapter file opens with an input of the shared preamble and carries no preamble text of its own, so the project's single Stacks-wide preamble staged at sources/stacks-duality/preamble.tex is excerpted verbatim under the recipe fixed by LANE\_BRIEF 22.1, as three spans, lines 8-13 (verbatim package and the reference/slogan/history environments), lines 14-18 (Xy-pic setup) and lines 154-217 (the native theorem declarations and mathematical macros). Nothing was added, removed or redefined; the wrapper class and the amsmath/amssymb/amsthm/graphicx base packages are supplied by the assembler.\par\medskip
\noindent\textit{Editorial note.} The counted statement and proof contain no cross reference, no citation and no included file, so the delivered body is self contained.\par\medskip
\noindent\textit{Editorial note.} Canonical location of the counted proposition: Stacks tag 0G4S (sources/stacks-duality/pinned-original-source.tex line 8320).\par\medskip
\noindent\textit{Allowed standard prerequisites (editorial).}\begin{itemize}
\item The Arin-Rees consequence identifying the pro-systems of cohomology. Same chapter, Lemma labelled lemma-consequence-Artin-Rees-bis, cited at proof lines 8343-8345. Gives the forward direction.
\item Truncation triangles. Derived Categories, Remark labelled derived-remark-truncation-distinguished-triangle, cited at proof lines 8349-8353. Supplies the compatible system of triangles for the induction.
\item The 2-out-of-3 property for Deligne systems. Same chapter, Lemma labelled lemma-deligne-system-2-out-of-3, cited at proof lines 8353-8355. Concludes the induction step.
\end{itemize}\medskip
\ifdefined\hypertarget\hypertarget{sourcepointer-1}{}\fi
\noindent\textit{Source reference.} Original lemma 0G4R: sources/stacks-duality/pinned-original-source.tex, source lines 8284--8284. Same-chapter pointer into the author environment lemma at source line 8284; invoked at proof line 8340 of the counted proof. It is retained only as a navigation pointer; every load-bearing step is inside the delivered span.\par\medskip
\ifdefined\hypertarget\hypertarget{sourcepointer-2}{}\fi
\noindent\textit{Source reference.} Original reference 08J5: sources/stacks-chow/tags.txt, source lines 10838--10838. Cross-chapter standard prerequisite located in the project tag list staged in the local snapshot; invoked at proof line 8351 of the counted proof. It is retained only as a navigation pointer; every load-bearing step is inside the delivered span.\par\medskip
\ifdefined\hypertarget\hypertarget{sourcepointer-3}{}\fi
\noindent\textit{Source reference.} Original lemma 0G4Q: sources/stacks-duality/pinned-original-source.tex, source lines 8236--8236. Same-chapter pointer into the author environment lemma at source line 8236; invoked at proof line 8354 of the counted proof. It is retained only as a navigation pointer; every load-bearing step is inside the delivered span.\par\medskip
\begin{lemma}
\label{lemma-characterize-extension-by-zero-algebra}
Let $j : U \to X$ be an open immersion of Noetherian schemes.
Let $a \leq b$ be integers. Let $(K_n)$ be an inverse system of
$D^b_{\textit{Coh}}(\mathcal{O}_X)$
such that $H^i(K_n) = 0$ for $i \not \in [a, b]$.
The following are equivalent
\begin{enumerate}
\item $(K_n)$ is pro-isomorphic to a Deligne system,
\item for every $p \in \mathbf{Z}$ there exists a coherent
$\mathcal{O}_X$-module $\mathcal{F}$ such that the pro-systems
$(H^p(K_n))$ and $(\mathcal{I}^n\mathcal{F})$ are pro-isomorphic.
\end{enumerate}
\end{lemma}

\begin{proof}
Assume (1). To prove (2) holds we may assume $(K_n)$ is a Deligne system.
By definition we may choose a bounded complex $\mathcal{F}^\bullet$
of coherent $\mathcal{O}_X$-modules and a quasi-coherent
sheaf of ideals cutting out $X \setminus U$ such that
$K_n$ is represented by $\mathcal{I}^n\mathcal{F}^\bullet$.
Thus the result follows from Lemma \ref{lemma-consequence-Artin-Rees-bis}.

\medskip\noindent
Assume (2). We will prove that $(K_n)$ is as in (1) by induction on
$b - a$. If $a = b$ then (1) holds essentially by assumption.
If $a < b$ then we consider the compatible system of
distinguished triangles
$$
\tau_{\leq a}K_n \to K_n \to \tau_{\geq a + 1}K_n \to (\tau_{\leq a}K_n)[1]
$$
See Derived Categories, Remark
\ref{derived-remark-truncation-distinguished-triangle}.
By induction on $b - a$ we know that $\tau_{\leq a}K_n$ and
$\tau_{\geq a + 1}K_n$ are pro-isomorphic to Deligne systems.
We conclude by Lemma \ref{lemma-deligne-system-2-out-of-3}.
\end{proof}

\end{document}
